\documentclass[11pt]{article}
\usepackage{setspace}
\usepackage[top=1in, bottom=1in, left=1in, right=1in]{geometry}
\usepackage{hyperref}
\usepackage{enumitem}
\usepackage[normalem]{ulem}
\usepackage{changepage}
\usepackage{xcolor}
\usepackage{ulem}

% =========================
% Custom macros for coauthor comments
% =========================
\newcommand{\sst}[1]{\textcolor{green}{#1}}
\newcommand{\tk}[1]{\textcolor{orange}{#1}}
\newcommand{\mw}[1]{\textcolor{magenta}{#1}}

\onehalfspacing

% Title information
\title{Note on Information Treatments}
\author{Tim Kircali, Stefanie Stantcheva, Matthias Weber}
\date{\today} % Use \date{} for no date

% Bibliography package
\usepackage[backend=biber, style=apa]{biblatex}
\addbibresource{references.bib} % Link to the bibliography file

\begin{document}
\maketitle
\section{Information Treatments}
\subsection{Evasion: Revenue Treatment}
\subsubsection*{General}
The first treatment provides respondents with information about the economic effects of third-party reporting and international information exchange in tax enforcement. We give them information about the size of the tax gap and we directly point out that third-party reporting and international exchange of inforamtion is effective in reducing tax evasion by explaining evidence which shows that tax evasion is rare in areas which are covered under third-party reporting and international exchange of information. The primary goal of the evasion:revenue treatment is to manipulate people’s knowledge about the economic effects of tax enforcement, third-party reporting and international information exchange without affecting considerations related to fairness \& inequality. \\ \\
Length of the video: 
\begin{itemize}
    \item New Version: 1:49
    \item Old Version: 3:07
\end{itemize}
\subsubsection*{Information Text: New Version}
The following explanations are used as subtitles and audio for an animated video.
\begin{itemize}
\item A major challenge for the US government is the loss of tax revenue through tax evasion.
\item Out of \$4.6 trillion in taxes, the IRS collects only 87\%. The missing 13\% add up to \$606 billion that never reach government revenue. 
\item The tax authorities have a number of options to combat tax evasion.
\item One option is to conduct tax audits, which involves closely examining tax returns.
\item Increasing tax audit activities is costly, but it helps the IRS to increase tax revenue.
\item Another powerful tool to combat tax evasion is third-party reporting. 
\item This means the IRS receives information directly from domestic firms and institutions.
\item Employers, for example, report their workers’ earnings to the IRS using Form W-2.
\item Another example is that banks use Form 1099 to report interest and dividend payments. 
\item Research shows that increasing third-party reporting reduces tax evasion and, therefore, increases tax revenue.
\item Beyond domestic reporting, international information exchange also plays a crucial role in reducing tax evasion. 
\item The IRS has information on most U.S. bank accounts, but some taxpayers hold accounts abroad.
\item Without international agreements, this information does not reach the IRS.
\item Research shows that expanding cooperation with more countries reduces tax evasion and, therefore, increases tax revenue.
\item To summarize the tools the IRS has to reduce tax evasion: It can increase tax audits, it can expand third-party reporting, and it can strengthen information exchange with other countries. All of this would increase tax revenue. 
\end{itemize}

\subsubsection*{Follow up questions}
We will ask 2-3 very basic questions to understand if respondents actually watched the video.
\subsection{Evasion: Inequality Treatment}
\subsubsection*{General}
The second specific information treatment again includes information on tax evasion, but also explains the fact that the tax gap is much more pronounced for high-income individuals and discusses why this is the case. After that, we discuss fairness considerations of tax evasion, such as the fact that honest taxpayers have to finance the fiscal gap caused by tax evaders. The primary goal of the evasion:inequality treatment is to manipulate people’s knowledge about the economic effects and distributional consequences of tax enforcement, third-party reporting and international information exchange.
\\ \\
Current length of the video: 
\begin{itemize}
    \item New Version: 2:31
    \item Old Version: 4:01
\end{itemize}
\subsubsection*{Information Text: New Version}
The following explanations are used as subtitles and audio for an animated video.
\begin{itemize}
\item A major challenge for the US government is the loss of tax revenue through tax evasion.
\item Out of \$4.6 trillion in taxes, the IRS collects only 87\%. The missing 13\% add up to \$606 billion that never reach government revenue. [\textit{Up to this point, the evasion: revenue and evasion: inequality treatments are identical.}]
\item Research has shown that not all groups of the population evade taxes to the same extent. 
\item While the poorer half of the population is responsible for only 11\% of evaded taxes, the 10\% richest taxpayers are responsible for 57\% 
\item But also within the wealthiest households, tax evasion is largely driven by a small group. 
\item The top 1\% richest taxpayers are responsible for every third dollar of evaded taxes.  
\item When some people don’t pay their taxes, the government either has to raise taxes or lower public services. Either way, all taxpayers are affected.
\item The tax authorities have a number of options to combat tax evasion. [\textit{From here on, the evasion: revenue and evasion: inequality treatments are identical.}]
\item The tax authorities have a number of options to combat tax evasion.
\item One option is to conduct tax audits, which involves closely examining tax returns.
\item Increasing tax audit activities is costly, but it helps the IRS to increase tax revenue.
\item Another powerful tool to combat tax evasion is third-party reporting. 
\item This means the IRS receives information directly from domestic firms and institutions.
\item Employers, for example, report their workers’ earnings to the IRS using Form W-2.
\item Another example is that banks use Form 1099 to report interest and dividend payments. 
\item Research shows that increasing third-party reporting reduces tax evasion and, therefore, increases tax revenue.
\item Beyond domestic reporting, international information exchange also plays a crucial role in reducing tax evasion. 
\item The IRS has information on most U.S. bank accounts, but some taxpayers hold accounts abroad.
\item Without international agreements, this information does not reach the IRS.
\item Research shows that expanding cooperation with more countries reduces tax evasion and, therefore, increases tax revenue.
\item To summarize the tools the IRS has to reduce tax evasion: It can increase tax audits, it can expand third-party reporting, and it can strengthen information exchange with other countries. All of this would increase tax revenue. 
\end{itemize}

\subsubsection*{Follow up questions}
We will ask 2-3 very basic questions to understand if respondents actually watched the video.
\subsection{Tax Filing Treatment}
\subsubsection*{General}
The third treatment discusses concerns about tax filing costs such as the fact that an average American taxpayer has a tax filing effort and cost of about 13 hours + \$250. It describes how (guided) tax software works and what a prepopulated tax return looks like. We emphasize that tax software and prepopulated tax returns help tax filers to make fewer mistakes and significantly reduce their tax filing costs. Importantly, we highlight the importance of third-party reporting to the tax administration in providing prepopulated tax returns. We explain that the IRS already has much tax-relevant data and that taxpayers would only have to give up a small share of data privacy to make it possible for the IRS to prepopulate tax returns. 
\\
\begin{itemize}
    \item New Version: 2:58
    \item Old Version: 3:26
\end{itemize}
\subsubsection*{Information Text: New Version}
The following explanations could be used as subtitles and audio for an animated video.
\begin{itemize}
    \item Most people need to file tax returns to report information to the tax administration. 
    \item This can be a tedious and expensive process.
    \item In the U.S., taxpayers spend an average of 13 hours to file their taxes.
    \item On average, additional financial costs of \$250 arise because many taxpayers either pay for advice, for someone preparing their tax returns, or for tax software. 
    \item Such software provides guidance through the tax formulas and explains which deductions might apply.
    \item It offers field-by-field advice, including clear explanations and examples. 
    \item In the U.S., most taxpayers pay commercial providers for tax software.
    \item However, the IRS could offer tax software free of charge, as tax authorities in many other countries already do.
    \item Not only could there be free tax software, but the IRS could also offer pre-populated tax returns. 
    \item This means that the tax returns are pre-filled with information the tax administration already has.
    \item To submit the tax return, some of the taxpayers would only need to enter missing information such as deductible expenses that the IRS does not know of. But many taxpayers would only have to briefly confirm that the pre-populated return is correct.
    \item This would save time, effort and it would reduce the potential of making mistakes when filing. 
    \item In order to offer pre-populated tax returns, the IRS needs the relevant data. Without these data, it will not be able to pre-populate many fields. 
    \item The tax authority often gets access to tax-relevant information by third-party reporting. 
    \item This means the IRS receives information directly from domestic firms and institutions. 
    \item Employers, for example, report their workers’ earnings to the IRS using Form W-2.  
    \item Another example is that banks use Form 1099 to report interest and dividend payments. 
    \item Currently, the IRS already receives enough third-party reported data to pre-populate around 45\% of tax returns. 
    \item However, these data are often sent and processed after the filing deadline.
    \item To pre-populate tax returns on time, the IRS needs to receive and process these data more quickly. 
    \item Another point is that other government institutions hold tax-related information, but these data are not shared with the IRS. 
    \item Local governments, for example, keep track of who owns a house and how much property tax they pay. However, they do not send this information to the IRS. 
    %\begin{itemize}
    %    \item \tk{Possibly we have to give here more detailed information about which data exactly should be exchanged.}
    %    \item \tk{Here we concentrate on data that is already in the hands of the government but not quick enough exchanged with the tax authority. Actually, also more data would help them but we decided to not concentrate on that. }
    %\end{itemize}
    \item Having access to these data would enable the IRS to pre-populate even more tax returns. 
    \item To summarize the most important steps towards pre-populating tax returns: 
    First, the speed of existing data transfer has to be increased. Second, additional tax-relevant data has to be reported by third parties.
\end{itemize}

\subsubsection*{Information Text}
\subsection{Full Info Treatment}
\subsubsection*{General}
The last information treatment includes information from all treatments. Essentially, it includes all the information from all the previous targeted videos. It first shows the tax filing treatment. After that, it transitions to the tax evasion: inequality treatment. There exists a slight overlap where third party reporting is explained in detail. This overlap is only shown in the first part of the video, but is not shown again in the second part.  
\\ \\
Current length of the video: 05:07
\subsubsection*{Information Text}
\begin{itemize}
    \item Most people need to file tax returns to report information to the tax administration. 
    \item This can be a tedious and expensive process.
    \item In the U.S., taxpayers spend an average of 13 hours to file their taxes.
    \item On average, additional financial costs of \$250 arise because many taxpayers either pay for advice, for someone preparing their tax returns, or for tax software. (separate financial cost slide)
    \item Such software provides guidance through the tax formulas and explains which deductions might apply.
    \item It offers field-by-field advice, including clear explanations and examples. 
    \item In the U.S., most taxpayers pay commercial providers for tax software.
    \item However, the IRS could offer tax software free of charge, as tax authorities in many other countries already do.
    \item Not only could there be free tax software, but the IRS could also offer pre-populated tax returns. 
    \item This means that the tax returns are pre-filled with information the tax administration already has.
    \item To submit the tax return, some of the taxpayers would only need to enter missing information such as deductible expenses that the IRS does not know of. But many taxpayers would only have to briefly confirm that the pre-populated return is correct.
    \item This would save time, effort and it would reduce the potential of making mistakes when filing. 
    \item In order to offer pre-populated tax returns, the IRS needs the relevant data. Without these data, it will not be able to pre-populate many fields. 
    \item The tax authority often gets access to tax-relevant information by third-party reporting. 
    \item This means the IRS receives information directly from domestic firms and institutions. 
    \item Employers, for example, report their workers’ earnings to the IRS using Form W-2.  
    \item Another example is that banks use Form 1099 to report interest and dividend payments. 
    \item Currently, the IRS already receives enough third-party reported data to pre-populate around 45\% of tax returns. 
    \item However, these data are often sent and processed after the filing deadline.
    \item To pre-populate tax returns on time, the IRS needs to receive and process these data more quickly. 
    \item Another point is that other government institutions hold tax-related information, but these data are not shared with the IRS. 
    \item Local governments, for example, keep track of who owns a house and how much property tax they pay. However, they do not send this information to the IRS. 
    \item Having access to these data would enable the IRS to pre-populate even more tax returns. 
    \item To summarize the most important steps towards pre-populating tax returns: 
    First, the speed of existing data transfer has to be increased. Second, additional tax-relevant data has to be reported by third parties. [\textit{Up to this point, the full filing treatment has been covered. Hereafter, the content of the Evasion: Inequality treatment will be shown. Small adjustments (which are noted) have been made to not cover content twice and to ensure a smooth tranisition.}]
    \item Third-party reporting would also help address another major challenge for the U.S. government: The loss of tax revenue due to evasion.  \textit{[Small adjustment to ensure a smooth transition]}
    \item Out of \$4.6 trillion in taxes, the IRS collects only 87\%. The missing 13\% add up to \$606 billion that never reach government revenue. [\textit{Up to this point of the evasion:inequality treatment, the evasion: revenue and evasion: inequality treatments are identical.}]
    \item Research has shown that not all groups of the population evade taxes to the same extent. 
    \item While the poorer half of the population is responsible for only 11\% of evaded taxes, the 10\% richest taxpayers are responsible for 57\% 
    \item But also within the wealthiest households, tax evasion is largely driven by a small group. 
    \item The top 1\% richest taxpayers are responsible for every third dollar of evaded taxes.  
    \item When some people don’t pay their taxes, the government either has to raise taxes or lower public services. Either way, all taxpayers are affected.
    \item The tax authorities have a number of options to combat tax evasion. [\textit{From here on, the evasion: revenue and evasion: inequality treatments are identical.}]
    \item The tax authorities have a number of options to combat tax evasion.
    \item One option is to conduct tax audits, which involves closely examining tax returns.
    \item Increasing tax audit activities is costly, but it helps the IRS to increase tax revenue.
    \item Another powerful tool to combat tax evasion is third-party reporting. 
    \item -[\textit{Content deleted due to overlap with filing treatment}]
    \item -[\textit{Content deleted due to overlap with filing treatment}]
    \item -[\textit{Content deleted due to overlap with filing treatment}]
    \item Research shows that increasing third-party reporting reduces tax evasion and, therefore, increases tax revenue.
    \item Beyond domestic reporting, international information exchange also plays a crucial role in reducing tax evasion. 
    \item The IRS has information on most U.S. bank accounts, but some taxpayers hold accounts abroad.
    \item Without international agreements, this information does not reach the IRS.
    \item Research shows that expanding cooperation with more countries reduces tax evasion and, therefore, increases tax revenue.
    \item To summarize the tools the IRS has to reduce tax evasion: It can increase tax audits, it can expand third-party reporting, and it can strengthen information exchange with other countries. All of this would increase tax revenue.
\end{itemize}
\subsubsection*{Follow up questions}
We will ask 2-3 very basic questions to understand if respondents actually watched all parts of the video.
\newpage
\section{Source of the information}
\subsubsection*{Tax Gap Estimation}
We use the following text:
\begin{itemize}
    \item "Out of \$4.6 trillion in taxes, the IRS collects only 87\%. The missing 13\% add up to \$606 billion that never reach government revenue."
\end{itemize}
These are tax gap projections for the year 2022 by \cite{IRS_TaxGap_2022_Pub5869}. The information used are the Total True Tax Liability (\$4,635) and the Net Tax Gap (Gross Tax Gap (\$696) - Enforced \& other late payments (\$90) = \$606). 
\subsubsection*{Distribution of Tax Evasion}
We use the following text:
\begin{itemize}
    \item "Research has shown that not all groups of the population evade taxes to the same extent." 
    \item "While the poorer half of the population cause only 11\% of tax evasion, the 10\% richest taxpayers are responsible for 57\%." 
    \item "But also within the wealthiest households, tax evasion is largely driven by a small group." 
    \item "The top 1\% richest taxpayers are responsible for every third dollar of evaded taxes. " 
    \item "Half of that amount is result of tax evasion by the super-rich, namely the top 0.1\% of the income distribution."
\end{itemize}
We use the information from \cite{guyton2021tax} (March 2021, Revised December 2023). Table 1: Share of Unreported Income. Column "Our Benchmark". The abstract of the study is: 
\vspace{1em}

\noindent We study tax evasion at the top of the U.S. income distribution using micro-data from random and operational audits and focused enforcement initiatives. Leveraging enforcement that revealed noncompliance ex post, we find that under the audit methods used during 2006–2013, individual random audit data failed to capture sophisticated evasion via offshore accounts and pass-through businesses. Consequently, estimates based solely on individual random audit data from this period under-state evasion by the highest-income Americans. We propose a theoretical explanation and construct new distributional estimates of noncompliance in the United States. Accounting for sophisticated evasion increases unreported income of the top 1\% of the income distribution in 2006–2013 by 50\% and increases the top 1\% fiscal income share by about 1 percentage point.
\vspace{0.5em} \\
\noindent In the following we can find some arguments why we we should use this data: 
\begin{itemize}
    \item Stratified random audits are commonly used to study and measure the extent of tax evasion. 
    \item The main difficulty is that general purpose audits may not detect all forms of tax evasion.
    \item The authors argue that these audits did not detect more sophisticated types of evasion. 
    \item They uncover two key limitations: tax evasion through foreign intermediaries (e.g., undeclared foreign bank accounts) and tax evasion via private businesses (e.g., under-reported revenues or overstated expenses in a partnership). 
    \item They analyse a sample of U.S. taxpayers who disclosed hidden offshore assets in the context of specific enforcement initiatives conducted in 2009-2012, and find 1) offshore tax evasion went almost entirely undetected in random audit data and 2) the very top of the income distribution is dramatically over-represented in this sample. 
    \item And they also find that tax evasion occurring in pass-through businesses (S corporations and partnerships) is highly under-detected in our individual random audit data. Pass-throughs are private businesses whose income is taxed as income to their owners; pass-through ownership and income are highly concentrated (Cooper et al., 2016).
    \item They then propose corrected estimates of underreported income through the income distribution. (which we use here)
\end{itemize}

\noindent Other relevant literature regarding the distribution of tax evasion can be found here: 
\begin{itemize}
    \item Distribution of tax evasion: \textcite{johns2010distribution}, \textcite{alstadsaeter2019tax}, \textcite{guyton2021tax}  
\end{itemize}
\subsubsection*{Tax Filing Costs}
We use the following text: 
\begin{itemize}
    \item "According to the IRS, US taxpayers need an average of 13 hours and \$250 to file their taxes."
\end{itemize}
This estimate of tax filing costs is provided by the IRS (\cite{IRS2022}). 
\subsubsection*{Prepopulated tax returns}
We provide the following text:
\begin{itemize}
    \item "Currently, the IRS already receives enough third-party reported data to prepopulate around 45\% of tax returns." 
\end{itemize}    
This estimate is the average of the two estimates \cite{goodman2023automated} provide. Using two different strategies, they provide the estimate that 42\% or 48\% of tax returns could be prepopulated (under the assumption that the IRS uses all this data).
\subsubsection*{Effectivity of Third Party Reporting}
We provide the following text: 
\begin{itemize}
    \item "Research has shown that tax evasion hardly takes place in areas that are covered by third-party reporting" 
    \item "Research has shown that having agreements on international information exchange is important to detect offshore tax evasion."
\end{itemize}
The relevant literature that supports these statements can be found here: 
\begin{itemize}
    \item Third-party reporting: \textcite{kleven2011unwilling}, \textcite{pomeranz2015no}, \textcite{kleven2016can}, \textcite{slemrod2017does}, \textcite{jensen2022employment} \\[0.1em]
    Automatic exchange of information: \textcite{baselgia2023compliance}, \textcite{boas2024taxing}
\end{itemize}
\section{Input Elevenlabs}
We use elevenlabs to generate audio. Here I documented the input that I used. 
\subsection{Evasion: Revenue}
Voice: Liam,
Model: Eleven v3,
Stability: Natural \\\\
Text: A major challenge for the US government is the loss of tax revenue through tax evasion. Out of 4.6 trillion dollars in taxes, the IRS collects only 87 percent. The missing 13 percent add up to [highlight]606 billion dollars that never reach government revenue. 
Research has shown that not all groups of the population evade taxes to the same extent. While the poorer half of the population is responsible for only 11 percent of evaded taxes, the 10 percent richest taxpayers are responsible for [highlight]57 percent. But also within the wealthiest households, tax evasion is largely driven by a small group. The top 1 percent richest taxpayers are responsible for every third dollar of evaded taxes.  
When some people don’t pay their taxes, the government either has to raise taxes or lower public services. Either way, [highlight: ALL]all taxpayers are affected.
The tax authorities have a number of options to combat tax evasion. One option is to conduct tax audits, which involves closely examining tax returns. Increasing tax audit activities is costly, but it helps the [not highlighting the word IRS]IRS to increase tax revenue.
Another powerful tool to combat tax evasion is third-party reporting. This means the [not highlighting the word IRS] IRS receives information directly from domestic firms and institutions. Employers, for example, report their workers’ earnings to the irs using Form W-2. Another example is that banks use Form 1099 to report interest and dividend payments. 
Research shows that increasing third-party reporting reduces tax evasion and, therefore, increases tax revenue.
Beyond domestic reporting, international information exchange also plays a crucial role in reducing tax evasion. The [not highlighting the word IRS]IRS has information on most U.S. bank accounts, but some taxpayers hold accounts abroad. Without international agreements, this information does not reach the [not highlighting the word IRS]IRS. Research shows that expanding cooperation with more countries reduces tax evasion and, therefore, increases tax revenue.
To summarize the tools the IRS has to reduce tax evasion: It can increase tax audits, it can expand third-party reporting, and it can strengthen information exchange with other countries. [Highlight]All of this would increase tax revenue. [not end here suddenly, it should fade out nicely!]
\subsection{Evasion: Revenue}
Voice: Liam,
Model: Eleven v3,
Stability: Natural\\\\
A major challenge for the US government is the loss of tax revenue through tax evasion. Out of 4.6 trillion dollars in taxes, the IRS collects only 87 percent. The missing 13 percent add up to [highlight]606 billion dollars that never reach government revenue. 
The tax authorities have a number of options to combat tax evasion. One option is to conduct tax audits, which involves closely examining tax returns. Increasing tax audit activities is costly, but it helps the [not highlighting the word IRS]IRS to increase tax revenue.
Another powerful tool to combat tax evasion is third-party reporting. This means the [not highlighting the word IRS] IRS receives information directly from domestic firms and institutions. Employers, for example, report their workers’ earnings to the irs using Form W-2. Another example is that banks use Form 1099 to report interest and dividend payments. 
Research shows that increasing third-party reporting reduces tax evasion and, therefore, increases tax revenue.
Beyond domestic reporting, international information exchange also plays a crucial role in reducing tax evasion. The [not highlighting the word IRS]IRS has information on most U.S. bank accounts, but some taxpayers hold accounts abroad. Without international agreements, this information does not reach the [not highlighting the word IRS]IRS. Research shows that expanding cooperation with more countries reduces tax evasion and, therefore, increases tax revenue.
To summarize the tools the IRS has to reduce tax evasion: It can increase tax audits, it can expand third-party reporting, and it can strengthen information exchange with other countries. [Highlight]All of this would increase tax revenue. [not end here suddenly, it should fade out nicely!]

\subsection{Filing Treatment}
Voice: Liam,
Model: Eleven v3,
Stability: Natural\\\\
Most people need to file tax returns to report information to the tax administration. This can be a tedious and expensive process. In the U.S., taxpayers spend an average of [highlight]13 hours to file their taxes. On average, additional financial costs of [highlight]250 dollars arise because many taxpayers either pay for advice, for someone preparing their tax returns, or for tax software. Such software provides guidance through the tax formulas and explains which deductions might apply. It offers field-by-field advice, including clear explanations and examples. 
In the U.S., most taxpayers pay commercial providers for tax software. However, the IRS could offer tax software [highlight]free of charge, as tax authorities in many other countries already do. Not only could there be free tax software, but the IRS could also offer pre-populated tax returns. This means that the tax returns are [highlight]pre-filled with information the tax administration already has. To submit the tax return, some of the taxpayers would only need to enter missing information such as deductible expenses that the IRS does not know of. But many taxpayers would only have to briefly confirm that the pre-populated return is correct. This would save time, effort and it would reduce the potential of making mistakes when filing.
In order to offer pre-populated tax returns, the IRS needs the relevant data. Without these data, it will not be able to pre-populate many fields. 
The tax authority often gets access to tax-relevant information by third-party reporting. This means the IRS receives information directly from domestic firms and institutions. Employers, for example, report their workers’ earnings to the IRS using Form W-2.  Another example is that banks use Form 1099 to report interest and dividend payments. 
Currently, [highlight]the IRS already receives enough third-party reported data to pre-populate around 45\% of tax returns. However, these data are often sent and processed after the filing deadline. To pre-populate tax returns on time, the IRS needs to receive and process these data more quickly. Another point is that other government institutions hold tax-related information, but these data are not shared with the IRS. Local governments, for example, keep track of who owns a house and how much property tax they pay. However, they do not send this information to the IRS. Having access to these data would enable the IRS to pre-populate even more tax returns. To summarize the most important steps towards pre-populating tax returns: 
First, the speed of existing data transfer has to be increased. Second, additional tax-relevant data has to be reported by third parties.[not end here suddenly, it should fade out nicely!]

\subsection{FilingxDataxLobbying}
Most people need to file tax returns to report information to the tax administration. This can be a tedious and expensive process. In the U.S., taxpayers spend an average of [highlight]13 hours to file their taxes. On average, additional financial costs of [highlight]250 dollars arise because many taxpayers either pay for advice, for someone preparing their tax returns, or for tax software. Such software provides guidance through the tax formulas and explains which deductions might apply. It offers field-by-field advice, including clear explanations and examples. 
In the U.S., most taxpayers pay commercial providers for tax software. However, the IRS could offer tax software [highlight]free of charge, as tax authorities in many other countries already do. Not only could there be free tax software, but the IRS could also offer pre-populated tax returns. This means that the tax returns are [highlight]pre-filled with information the tax administration already has. To submit the tax return, some of the taxpayers would only need to enter missing information such as deductible expenses that the IRS does not know of. But many taxpayers would only have to briefly confirm that the pre-populated return is correct. This would save time, effort and it would reduce the potential of making mistakes when filing.
In order to offer pre-populated tax returns, the IRS needs the relevant data. Without these data, it will not be able to pre-populate many fields. 
The tax authority often gets access to tax-relevant information by third-party reporting. This means the IRS receives information directly from domestic firms and institutions. Employers, for example, report their workers’ earnings to the IRS using Form W-2.  Another example is that banks use Form 1099 to report interest and dividend payments. 
Currently, [highlight]the IRS already receives enough third-party reported data to pre-populate around 45\% of tax returns. However, these data are often sent and processed after the filing deadline. To pre-populate tax returns on time, the IRS needs to receive and process these data more quickly. Another point is that other government institutions hold tax-related information, but these data are not shared with the IRS. Local governments, for example, keep track of who owns a house and how much property tax they pay. However, they do not send this information to the IRS. Having access to these data would enable the IRS to pre-populate even more tax returns. [Pause] Limited data access is not the only reason the IRS has not yet introduced better filing services. Such reforms also require support from federal politicians. And when it comes to policy debates about better filing services, the tax software industry is very active in putting forward its position. Over the last 20 years, two providers alone have spent around 100 million dollars on lobbying. These companies have a strong financial interest in maintaining the current system. Their tax preparation business generates billions in revenue each year. Simpler filing would reduce demand for their services. As long as they are able to persuade politicians to maintain the current system, reforms will be difficult to implement.[not end here suddenly, it should fade out nicely!]

\subsection{FilingxData}
Most people need to file tax returns to report information to the tax administration. This can be a tedious and expensive process. In the U.S., taxpayers spend an average of [highlight]13 hours to file their taxes. On average, additional financial costs of [highlight]250 dollars arise because many taxpayers either pay for advice, for someone preparing their tax returns, or for tax software. Such software provides guidance through the tax formulas and explains which deductions might apply. It offers field-by-field advice, including clear explanations and examples. 
In the U.S., most taxpayers pay commercial providers for tax software. However, the IRS could offer tax software [highlight]free of charge, as tax authorities in many other countries already do. Not only could there be free tax software, but the IRS could also offer pre-populated tax returns. This means that the tax returns are [highlight]pre-filled with information the tax administration already has. To submit the tax return, some of the taxpayers would only need to enter missing information such as deductible expenses that the IRS does not know of. But many taxpayers would only have to briefly confirm that the pre-populated return is correct. This would save time, effort and it would reduce the potential of making mistakes when filing.
In order to offer pre-populated tax returns, the IRS needs the relevant data. Without these data, it will not be able to pre-populate many fields. 
The tax authority often gets access to tax-relevant information by third-party reporting. This means the IRS receives information directly from domestic firms and institutions. Employers, for example, report their workers’ earnings to the IRS using Form W-2.  Another example is that banks use Form 1099 to report interest and dividend payments. 
Currently, [highlight]the IRS already receives enough third-party reported data to pre-populate around 45\% of tax returns. However, these data are often sent and processed after the filing deadline. To pre-populate tax returns on time, the IRS needs to receive and process these data more quickly. Another point is that other government institutions hold tax-related information, but these data are not shared with the IRS. Local governments, for example, keep track of who owns a house and how much property tax they pay. However, they do not send this information to the IRS. Having access to these data would enable the IRS to pre-populate even more tax returns. 

\subsection{FilingxLobbying}
Most people need to file tax returns to report information to the tax administration. This can be a tedious and expensive process. In the U.S., taxpayers spend an average of [highlight]13 hours to file their taxes. On average, additional financial costs of [highlight]250 dollars arise because many taxpayers either pay for advice, for someone preparing their tax returns, or for tax software. Such software provides guidance through the tax formulas and explains which deductions might apply. It offers field-by-field advice, including clear explanations and examples. 
In the U.S., most taxpayers pay commercial providers for tax software. However, the IRS could offer tax software [highlight]free of charge, as tax authorities in many other countries already do. Not only could there be free tax software, but the IRS could also offer pre-populated tax returns. This means that the tax returns are [highlight]pre-filled with information the tax administration already has. To submit the tax return, some of the taxpayers would only need to enter missing information such as deductible expenses that the IRS does not know of. But many taxpayers would only have to briefly confirm that the pre-populated return is correct. This would save time, effort and it would reduce the potential of making mistakes when filing.
Reforms such as offering free tax software or prepopulated tax returns requires support from federal politicians. And when it comes to policy debates about better filing services, the tax software industry is very active in putting forward its position. Over the last 20 years, two providers alone have spent around 100 million dollars on lobbying. These companies have a strong financial interest in maintaining the current system. Their tax preparation business generates billions in revenue each year. Simpler filing would reduce demand for their services. As long as they are able to persuade politicians to maintain the current system, reforms will be difficult to implement.[not end here suddenly, it should fade out nicely!]


\subsection{Full info}
Voice: Liam
Model: Eleven v3
Stability: Natural\\\\
Most people need to file tax returns to report information to the tax administration. This can be a tedious and expensive process. In the U.S., taxpayers spend an average of [highlight]13 hours to file their taxes. On average, additional financial costs of [highlight]250 dollars arise because many taxpayers either pay for advice, for someone preparing their tax returns, or for tax software. Such software provides guidance through the tax formulas and explains which deductions might apply. It offers field-by-field advice, including clear explanations and examples. In the U.S., most taxpayers pay commercial providers for tax software. However, the IRS could offer tax software free of charge, as tax authorities in many other countries already do. Not only could there be free tax software, but the IRS could also offer pre-populated tax returns. This means that the tax returns are [highlight]pre-filled with information the tax administration already has. To submit the tax return, some of the taxpayers would only need to enter missing information such as deductible expenses that the IRS does not know of. But many taxpayers would only have to briefly confirm that the pre-populated return is correct. This would save time, effort and it would reduce the potential of making mistakes when filing. In order to offer pre-populated tax returns, the IRS needs the relevant data. Without these data, it will not be able to pre-populate many fields. The tax authority often gets access to tax-relevant information by third-party reporting. This means the IRS receives information directly from domestic firms and institutions. Employers, for example, report their workers’ earnings to the IRS using Form W-2. Another example is that banks use Form 1099 to report interest and dividend payments. Currently, [highlight]the IRS already receives enough third-party reported data to pre-populate around 45 percent of tax returns. However, these data are often sent and processed after the filing deadline. To pre-populate tax returns on time, the IRS needs to receive and process these data more quickly. Another point is that other government institutions hold tax-related information, but these data are not shared with the IRS. Local governments, for example, keep track of who owns a house and how much property tax they pay. However, they do not send this information to the IRS. Having access to these data would enable the IRS to pre-populate even more tax returns. To summarize the most important steps towards pre-populating tax returns: First, the speed of existing data transfer has to be increased. Second, additional tax-relevant data has to be reported by third parties. [short pause] Third-party reporting would also help address another major challenge for the U.S. government: the loss of tax revenue due to evasion. Out of 4.6 trillion dollars in taxes, the IRS collects only 87 percent. The missing 13 persent add up to 606 billion dollars that never reach government revenue. Research has shown that not all groups of the population evade taxes to the same extent. While the poorer half of the population is responsible for only 11 percent of evaded taxes, the 10 percent richest taxpayers are responsible for [highlight]57 percent. But also within the wealthiest households, tax evasion is largely driven by a small group. The top 1 percent richest taxpayers are responsible for every third dollar of evaded taxes.  When some people don’t pay their taxes, the government either has to raise taxes or lower public services. Either way, [highlight: ALL]all taxpayers are affected. The tax authorities have a number of options to combat tax evasion. One option is to conduct tax audits, which involves closely examining tax returns. Increasing tax audit activities is costly, but it helps the IRS to increase tax revenue. Another powerful tool to combat tax evasion is third-party reporting. Research shows that increasing third-party reporting reduces tax evasion and, therefore, increases tax revenue. Beyond domestic reporting, international information exchange also plays a crucial role in reducing tax evasion. The IRS has information on most U.S. bank accounts, but some taxpayers hold accounts abroad. Without international agreements, this information does not reach the IRS. Research shows that expanding cooperation with more countries reduces tax evasion and, therefore, increases tax revenue. To summarize the tools the IRS has to reduce tax evasion: It can increase tax audits, it can expand third-party reporting, and it can strengthen information exchange with other countries. [Highlight]All of this would increase tax revenue. [not end here suddenly, it should fade out nicely!]
\newpage
\printbibliography

\end{document}